% 4.6 章 · 谱方法与伪谱
% 这一章要把"用谱方法解 PDE"拆成四件互不相同的数学事实：
%   1) 谱微分是一个代数恒等式（三角多项式上是恒等算子），不是"精度很高的近似"；
%   2) 收敛速度由光滑性决定：解析 $\Rightarrow$ 指数，C^0 尖角 $\Rightarrow$ N^-1，跳变 $\Rightarrow$ 不收敛（Gibbs）；
%   3) 伪谱乘积的误差必须拆成三笔账：混叠、2/3 截断、补零——它们量级不同、代价不同；
%   4) 时间推进的去混叠只改非线性项，不改格式的阶；"能量守恒改善"不是它的卖点。
% 所有数字都来自本仓库 tests/4.6/ 的运行结果；断言附表由 tools/build_book.py 生成，
% \input 进来，不手抄。

\numlabchapter{谱方法与伪谱}{4.6}

\section{这一章要解决什么}

给定周期函数 $u$ 在 $N$ 个等距节点上的值，谱方法把解写成有限三角级数
\[
u_N(x)=\sum_{|k|\le N/2}\hat u_k e^{2\pi i kx/L},\qquad
\hat u_k=\frac1N\sum_{j=0}^{N-1}u(x_j)e^{-2\pi i kx_j/L},
\]
于是"微分"变成逐模态乘法、"非线性乘积"变成逐点乘法，而后者正是所有麻烦的来源。
这一章回答四个问题，每个都配上实测数字：

\begin{enumerate}
  \item \textbf{谱微分到底有多准？} 在三角多项式上它是\textbf{恒等算子}：
        $u=\sin 3x+\tfrac12\cos 7x$ 时，$N=16$ 与 $N=64$ 的相对误差是
        $4.66\times10^{-15}$ 与 $1.01\times10^{-14}$——与 $N$ 无关，就是舍入。
        同一个算子在锯齿剖面上让 $\max|u'|$ 随 $N$ 线性涨：$6.73,13.80,27.92$
        （$N=32,64,128$）。所以"谱方法准不准"的正确答案是"看数据光滑不光滑"。
  \item \textbf{收敛速度从哪来？} 模态系数由光滑性定阶：跳变 $k^{-1}$、$C^0$ 尖角 $k^{-2}$、
        多项式凸包 $k^{-4}$（实测 $-0.9921,-1.9840,-3.9999$，窗口 $k\in[8,48]$，$R^2\ge0.9999$）；
        逼近误差再快一阶（尖角实测指数 $+1.0038$），而间断数据的误差\textbf{不下降}
        （$N=16\to256$ 时 $1.9078\to1.8022$，比值 $1.06$）。Gibbs 过冲收敛到
        $\frac2\pi\mathrm{Si}(\pi)-1=0.1789797$，不是 $0$。
  \item \textbf{伪谱乘积错在哪？} 三元组 $p+q=\pm N$ 的折返（混叠）铺满整个保留带，
        实测相对偏差 $1.196\times10^{-2}$；而 $2/3$ 截断自己要丢掉 $3.740\times10^{-2}$。
        两项必须分开报：截断的代价比它顺手消掉的混叠还大 $3$ 倍。
  \item \textbf{去混叠换来了什么？} 对已经 $2/3$ 截断的场，去混叠乘积与朴素乘积的
        能量恒等式残差都是舍入级（$3.4\times10^{-18}$ 与 $8.4\times10^{-17}$），
        所以"能量守恒"不是它的证据；它修的是非线性项在保留带里的取值。
        在分辨良好的 Burgers 算例里，两种算法的总能量相对差只有 $2\times10^{-4}$，
        差别出现在 Enstrophy（$3.5\%$）与逐点解（$\max|\Delta u|/|u|=3.9\times10^{-2}$）上。
\end{enumerate}

记号：$N$ 是节点数（自由度），$L$ 是周期长度，$\kappa_k=2\pi k/L$ 是波数，
$k_d=1/\sqrt{\nu t}$ 是黏性截断波数，$\varepsilon=2\nu\sum_k\kappa_k^2|\hat u_k|^2$
是耗散率，$\Pi(k)$ 是能流。本章的"实测"数字全部取自 \texttt{tests/4.6/} 的断言记录
（共 $74$ 条：A 类 $23$、B 类 $33$、C 类 $18$）。

\section{谱微分：一个恒等式，以及它的边界}

\begin{theorem}[谱微分]
设 $u(x)=\sum_{|k|\le K}\hat u_k e^{i\kappa_k x}$ 是带宽 $K\le N/2$ 的三角多项式，则
$\widehat{u'}_k=i\kappa_k\hat u_k$ 对每个 $|k|\le K$ 精确成立，且该式与 $N$ 无关。
\end{theorem}

\begin{proof}
逐项求导：$\frac{d}{dx}e^{i\kappa_k x}=i\kappa_k e^{i\kappa_k x}$，有限和可以逐项微分。
由于 $K\le N/2$，每个模态都在 $N$ 点网格的表示范围内（$k$ 与 $k+N$ 不混淆），
所以 DFT 系数就是 $\hat u_k$ 本身（无混叠）。于是离散微分算子作用在系数上
只是乘以 $i\kappa_k$，误差只来自浮点乘加。
\end{proof}

\begin{corollary}[误差是 $u$ 的量级，不是"很小"]
实测：$N=16$ 时 $4.66\times10^{-15}$，$N=64$ 时 $1.01\times10^{-14}$（相对 $\max|u'|$）。
两个数都在 $10^{-14}$ 附近而不是 $10^{-2}$ 附近，说明实现没有额外的近似；
但也不随 $N$ 变小——因为已经没有可减的误差了。
\end{corollary}

\begin{theorem}[间断数据的谱微分不收敛]
设 $u$ 有跳变。则存在常数 $c>0$ 使得 $\|u'_N\|_\infty\ge cN$：Gibbs 层的宽度是
$O(1/N)$、幅度是 $O(1)$，导数幅度因此是 $O(N)$。
\end{theorem}

\begin{proof}
跳变函数的截断级数在跳变两侧出现幅度约 $0.0895$ 倍跳变量的过冲，振荡周期 $\sim L/N$。
把振荡写成 $A\,g(Nx/L)$（$g$ 是固定波形），$u'_N\sim A(N/L)g'$，故 $\|u'_N\|_\infty=O(N)$。
\end{proof}

\textbf{实测}：锯齿剖面上 $\max|u'|=6.73,13.80,27.92$（$N=32,64,128$），
增长 $4.15$ 倍（$N$ 增长 $4$ 倍）。所以"谱方法总能微分"是错的：
判据是数据光滑，不是方法。

\section{系数衰减、逼近误差与 Gibbs}

\begin{theorem}[代数衰减的阶]
若 $u$ 是 $C^{p-1}$ 且第 $p$ 阶导数有跳变，则 $\hat u_k=O(k^{-(p+1)})$。
特别地：跳变（$p=1$）给 $k^{-1}$，$C^0$ 尖角（$p=2$）给 $k^{-2}$。
\end{theorem}

\begin{proof}
对 $\hat u_k$ 分部积分 $p-1$ 次，把导数转移到 $u$ 上：
$\hat u_k=(i\kappa_k)^{-(p-1)}\widehat{u^{(p-1)}}_k$。对 $C^0$ 且分段 $C^1$ 的函数
（第 $p$ 阶导数有跳变 $\Delta$），$\widehat{u^{(p-1)}}_k\to\Delta/(2\pi i k)$，
故 $\hat u_k\sim \Delta/(i\kappa_k)^p$，即 $O(k^{-p})$；再写一遍分部积分得到 $O(k^{-(p+1)})$
的上界形式。多项式凸包（$C^3$、四阶导数跳变）给 $k^{-4}$。
\end{proof}

\textbf{实测}（$n=512$，窗口 $k\in[8,48]$）：
锯齿 $-0.9921$（$R^2=0.999988$）、尖顶 $-1.9840$（$R^2=0.999989$）、
多项式凸包 $-3.9999$（$R^2=1.000000$）。解析剖面（Poisson 核）在同一框架下
\textbf{没有有限的代数指数}：分段平均斜率 $-3.06,-6.64,-12.77$（$k\in[2,8],[8,16],[16,32]$），
每段至少再陡 $3.6$。

\begin{corollary}[逼近误差比系数快一阶]
$\|u-S_Nu\|_\infty\le\sum_{|k|>N/2}|\hat u_k|=O(N^{-p})$：$C^0$ 尖角的误差按 $N^{-1}$ 下降。
\end{corollary}

\begin{proof}
误差是尾部之和；用 $|\hat u_k|\le Ck^{-(p+1)}$ 与积分判别法，
$\sum_{k>N/2}k^{-(p+1)}=O(N^{-p})$。
\end{proof}

\textbf{实测}：尖角误差 $0.024,0.012,0.006,0.003,0.001$（$N=16,32,64,128,256$），
拟合指数 $+1.0038$（$R^2=0.999990$）。跳变那条：$N=16$ 时 $1.9078$、$N=256$ 时 $1.8022$，
比值 $1.06$——\textbf{不下降}。解析剖面在 $N=256$ 时 $7.1\times10^{-15}$（撞地板）。

\begin{theorem}[Gibbs 过冲的极限]
方波的 $K$ 项部分和在跳变附近的过冲收敛到
$\frac2\pi\mathrm{Si}(\pi)-1=0.1789797$，其中 $\mathrm{Si}(\pi)=\int_0^\pi\frac{\sin t}{t}dt$；
即过冲的\textbf{幅度}有非零极限，$K$ 增大只压窄振荡。
\end{theorem}

\begin{proof}
部分和的误差是 Dirichlet 核的卷积；在跳变右侧 $x=\pi/(2K+1)$ 附近取极值，
误差化为 $\frac1\pi\int_0^\pi\frac{\sin t}{t}dt-\frac12$ 的常数部分。
\end{proof}

\textbf{实测}：$K=4,64,1024$ 时 $0.182327,0.178918,0.178391$，
与极限的差 $5.89\times10^{-4}$，且单调下降。

\begin{theorem}[指数收敛的可执行判据]
若存在 $\rho<1$ 使 $|\hat u_k|\le C\rho^{|k|}$，则每把 $N$ 加倍，逼近误差至少再乘一个
\textbf{与 $N$ 无关且 $<1$ 的因子}；反之，若 $\hat u_k\sim Ck^{-p}$，则每加倍误差比恒为
$2^{-p}$（一个常数）。因此判据是"每加倍比值是否\textbf{递减}"，而不是某个拟合出的 $p$。
\end{theorem}

\begin{proof}
尾部 $\sum_{k>N/2}\rho^k\propto\rho^{N/2}$，$N\to2N$ 时比值为 $\rho^{N/2}$，
它随 $N$ 指数变小且始终 $<1$（所以比值序列递减、没有下界）。
代数情形 $\sum_{k>N/2}k^{-(p+1)}\propto N^{-p}$，比值恒为 $2^{-p}$。
\end{proof}

\textbf{实测}：默认帧（$N=64,\rho=0.6$）的每加倍比值 $0.143\to4.08\times10^{-4}$；
平流算例的谱误差比值 $5.55\times10^{-2},1.20\times10^{-2},1.23\times10^{-3},4.07\times10^{-5},2.90\times10^{-8}$
（无下界）；同算例二阶差分是 $0.657,0.5413,0.3352,0.2548$（有下界 $1/4$）。
拟合窗口 $[16,32],[24,48],[32,64]$ 上的谱斜率 $5.03,8.39,11.57$ 逐段变陡。

\begin{remark}[拟合出的 $p$ 会骗人]
默认帧对 $|a_k|$ 做双对数最小二乘，窗口 $k\in[4,32]$ 给 $p=-7.129$、$R^2=0.938$——
比任何代数衰减都陡。但同一曲线的局部斜率从 $-0.737$ 走到 $-3.942$；
把它当"阶"用就等于把窗口当成结论。
\end{remark}

\section{Chebyshev：换基不换结论}

\begin{theorem}[Chebyshev 微分矩阵对多项式精确]
在 $n+1$ 个 Gauss--Lobatto 点 $x_j=\cos(\pi j/n)$ 上，Trefethen 意义下的微分矩阵
$D$ 作用在次数 $\le n$ 的多项式上给出\textbf{精确}导数。
\end{theorem}

\begin{proof}
$D_{ij}$ 由 Lagrange 基多项式的导数闭式给出；对任意次数 $\le n$ 的多项式 $p$，
$p'$ 是它自己的 $n$ 次插值多项式的导数，而 $(Dp)_i$ 恰是该插值多项式在 $x_i$ 的导数。
\end{proof}

\textbf{实测}：$n=16$ 时 $D x$ 与 $1$ 的最大偏差 $3.6\times10^{-15}$，$D x^2$ 与 $2x$ 是
$1.4\times10^{-14}$；$n=32$ 时 $5.7\times10^{-14}$ 与 $7.1\times10^{-14}$。
矩阵行和 $2.8\times10^{-14}$（常数导数为零）。端点 $x_0=1,x_n=-1$ 精确。

\begin{theorem}[系数衰减与误差：与 Fourier 侧同构]
解析函数（Runge 函数）的 Chebyshev 系数按 $\rho^k$ 指数衰减，$\rho=3-2\sqrt2=0.1716$；
$C^0$ 尖角（$|x|$）按 $k^{-2}$；跳变（$\mathrm{sign}\,x$）按 $k^{-1}$。
\end{theorem}

\begin{proof}
Chebyshev 系数是 $f(\cos\theta)$ 的 Fourier 余弦系数；Runge 函数 $1/(1+25x^2)$ 的
奇点在 $x=\pm i/5$，映射到 Bernstein 椭圆参数 $\rho=3-2\sqrt2$。$|x|$ 与 $\mathrm{sign}\,x$
用上一节的分部积分结论（在 $\theta$ 坐标下同样的跳变结构）。
\end{proof}

\textbf{实测}：Runge 的误差每加倍比值 $0.0765,0.0202,0.0045,4.8\times10^{-4},1.9\times10^{-5},2.2\times10^{-7}$
（严格递减）；$|x|$ 的误差指数 $+0.9990$（$n=16..128$）；$\mathrm{sign}$ 的误差
$1.9944\to1.9540$（比值 $1.02$，不下降）。

\begin{corollary}[非光滑数据换基也没用]
$|x|$ 的导数误差在 $n=16,32,64,128$ 上是 $2.3881,2.3867,2.3864,2.3863$：
$O(1)$ 且与 $n$ 无关。Chebyshev 的红利是"少数点达到同样的分辨率"，不是绕过光滑性。
\end{corollary}

\begin{proof}
$|x|$ 的导数在 $x=0$ 有幅度 $2$ 的跳变；插值多项式的导数在该点附近同样出现
$O(1)$ 的过冲，与 $n$ 无关（与上一节的 Gibbs 现象同源，只是权重不同）。
\end{proof}

\begin{remark}[零点模态会伪造斜率]
采样方波的偶次谐波不是 $0$ 而是常数 $1/N$（$N=64$ 时 $0.03125$），
拟合被拖成 $p=-0.848,R^2=0.346$；Chebyshev 的 Runge 是偶函数、奇次系数恰好为 $0$
（实测 $2.7\times10^{-16}$），不筛就得到 $+56$ 这种正斜率。
$n=64$ 的 Chebyshev 默认帧筛掉了 $14$ 个零模态后才得到 $p=-2.749,R^2=0.927$。
\end{remark}

\section{伪谱：三笔账必须分开算}

\begin{theorem}[混叠的落点]
$N$ 点网格上的逐点乘积等价于离散谱的\textbf{循环卷积}：乘积谱中模态 $p+q$ 若
$|p+q|>N/2$，就被折到 $p+q\mp N$。折返落点铺满整个保留带，因此混叠不只污染高波数。
\end{theorem}

\begin{proof}
$N$ 点 DFT 的基满足 $e^{2\pi i(p+q)x_j/N}=e^{2\pi i(p+q\mp N)x_j/N}$，
即波数在模 $N$ 意义下等价；逐点乘积的变换因此是循环卷积。
\end{proof}

\textbf{实测（手算可验的现场）}：$N=12$、$u=\cos(2\pi\cdot4x)$ 时 $u^2=0.5+0.5\cos(2\pi\cdot8x)$，
模态 $8$ 折回 $8-12=-4$，于是朴素乘积在 $k=-4$ 处出现 $0.25$，
而高分辨率参照与去混叠乘积在该处都是 $0$（$<10^{-14}$）。

\begin{theorem}[$3N/2$ 补零的保证]
设输入带宽 $K$，补零到 $M\ge3K$ 后相乘，则乘积在保留带 $|k|<M/2$ 内无混叠。
特别地：全带宽输入（$K=N/2$）配 $M=3N/2$ 时，保留带内唯一的污染是
$|k|=N/2$ 那条边界模态（来自 $p=q=N/2$ 的自乘积）。
\end{theorem}

\begin{proof}
乘积的真实带宽是 $2K$；在 $M$ 点网格上，模 $>M/2$ 的部分折回后的最小落点是
$|2K-M|$（取 $K$ 最大时）或更一般地 $M-2K$ 的镜像。要让它落在保留带
$|k|<K$ 之外，需要 $M-2K\ge K$，即 $M\ge3K$。取 $K=N/2$、$M=3N/2$ 得
$M-2K=N/2$，恰好只有边界模态被覆盖。
\end{proof}

\textbf{实测}（$N=128$、$|a_k|=k^{-1}$、随机相位、种子固定）：
只补零到精确 $M=3N/2$、完全不截断时，相对全带参照的偏差是 $7.76\times10^{-5}$，
而且残差全部落在 $|k|=N/2$（低带偏差 $6\times10^{-7}$ 与 Nyquist 偏差相等即证据）；
本内核默认把 $M$ 补到不小于 $3N/2$ 的 $2$ 的幂（$N=128$ 时 $M=256=2N$），偏差 $8\times10^{-18}$。

\begin{theorem}[$2/3$ 规则是另一条路]
把输入截到 $|k|\le N/3$ 后，$N$ 点网格上的朴素乘积在保留带 $|k|\le N/3$ 内也是精确的：
折返落点是 $|2N/3-N|=N/3$，正好在带外。代价是丢掉 $|k|>N/3$ 的输入能量。
\end{theorem}

\begin{proof}
带宽 $K=N/3$ 的乘积带宽 $2N/3>N/2$，折返把 $k\in(N/2,2N/3]$ 映到
$k-N\in(-N/2,-N/3]$，即 $|k|\in[N/3,N/2)$，都在保留带之外。
\end{proof}

\textbf{实测}：$N=128$、$|a_k|=k^{-1}$ 时，$2/3$ 截断丢掉的输入能量占比 $5.37\times10^{-3}$
（$\mathrm{decay}=0.5$ 时 $8.80\times10^{-2}$），乘积谱的整体偏差 $3.740\times10^{-2}$；
而"$2/3$ 截断 $+$ $3/2$ 补零"对截断后的场是\textbf{逐位精确}的（偏差 $0$，贴 $10^{-16}$）。

\begin{center}
\begin{tabular}{@{}lrr@{}}
\toprule
量（$N=128$、$|a_k|=k^{-\mathrm{decay}}$） & $\mathrm{decay}=1$ & $\mathrm{decay}=0.5$ \\
\midrule
混叠误差（朴素 vs 全带参照） & $1.196\times10^{-2}$ & $6.64\times10^{-2}$ \\
其中落在保留带 $|k|\le N/3$ & $2.67\times10^{-3}$ & $3.78\times10^{-2}$ \\
$2/3$ 截断代价（截断 vs 全带参照） & $3.740\times10^{-2}$ & $1.40\times10^{-1}$ \\
只补零（$M=3N/2$、不截断） & $7.76\times10^{-5}$ & $1.60\times10^{-3}$ \\
两条路线总差（朴素 vs $2/3+3/2$） & $4.034\times10^{-2}$ & $1.53\times10^{-1}$ \\
\bottomrule
\end{tabular}
\end{center}

\begin{corollary}[截断的代价大于它消掉的混叠]
上表第 1 行与第 3 行之比是 $3.13$（$\mathrm{decay}=1$）；把 $\mathrm{decay}$ 扫过
$0.5,1,2,2.5$，比值稳定在 $2.10,3.13,5.65,6.74$。所以"去混叠改善了精度"这句话
必须说清报的是哪一笔账。
\end{corollary}

\begin{remark}[能量恒等式残差不是判据]
$R=\sum_k\kappa_k\,\mathrm{Im}(\overline{\hat u_k}\hat w_k)$ 对已 $2/3$ 截断的场，
朴素乘积与去混叠乘积分别是 $3.4\times10^{-18}$ 与 $8.4\times10^{-17}$（都看不到混叠，
因为折返落在保留带之外）；把场放宽到全带宽，$R$ 变成 $10^{-3}\sim10^{-1}$ 量级，
而且\textbf{同一张谱只换相位}就能从 $7.65\times10^{-18}$（二次相位）跳到
$4.05\times10^{-3}$（线性相位）、$7.17\times10^{-3}$（随机相位）——跨度 $9.4\times10^{14}$。
判据只能是"非线性项自身的偏差"。
\end{remark}

\section{时间推进：积分因子 RK4}

把 $u_t+uu_x=\nu u_{xx}$ 写成系数空间的常微分方程组
\[
\frac{d\hat u_k}{dt}=\underbrace{-\frac{i\kappa_k}{2}\widehat{(u^2)}_k}_{\text{非线性}}
-\underbrace{\nu\kappa_k^2\hat u_k}_{\text{线性}},
\]
把线性部分精确积分出来（积分因子 $e^{-\nu\kappa_k^2t}$），非线性部分用经典 RK4。

\begin{theorem}[阶与稳定界]
积分因子 RK4 在每个固定 $\kappa_k$ 上是四阶的；线性项的指数因子把黏性稳定限制
$dt<2.78/(\nu\kappa_{\max}^2)$ \textbf{整个消掉}，剩下的是对流 CFL 限制
$dt\le c\,dx/\max|u|$（本内核 $c=0.35$）。
\end{theorem}

\begin{proof}
指数因子是线性方程的精确解算子，与 RK4 复合不降低非线性部分的四阶精度
（线性部分不再引入截断误差）；黏性稳定界来自显式 RK4 的稳定多项式在
$z=-\nu\kappa^2dt$ 上的边界，改用精确指数后该约束不存在。
\end{proof}

\textbf{实测}：对 $N=64,\nu=0.01$ 的算例做自收敛（基准 $dt=T/640$），
实测时间阶 $3.9932$（$R^2=0.999999$），每减半步长的误差比
$0.0632,0.0627,0.0625$（理论 $1/16=0.0625$）；与 \texttt{DOP853}
（$rtol=10^{-12}$，同一条半离散方程）的相对偏差在 $dt=0.0025$ 时是 $1.43\times10^{-12}$，
在 $dt=0.04$ 时是 $9.30\times10^{-8}$（相差 $6.5\times10^{4}$ 倍）。

\section{能流与能量收支}

\begin{theorem}[逐带能量恒等式]
记 $E_k=|\hat u_k|^2$、$E_{\le k}=\sum_{|j|\le k}E_j$，则
$\Pi(k)=-\big(\dot E_{\le k}+2\nu\sum_{|j|\le k}\kappa_j^2E_j\big)$，
其中 $\Pi(k)=\sum_{|j|\le k}\kappa_j\,\mathrm{Im}(\overline{\hat u_j}\hat w_j)$
是 $k$ 处的能流（$\hat w=\widehat{(u^2)}$）。
\end{theorem}

\begin{proof}
对截断的 Galerkin 系统 $\dot{\hat u}_j=-i\kappa_j\hat w_j/2-\nu\kappa_j^2\hat u_j$
取模平方、对 $|j|\le k$ 求和，非线性项的贡献正是 $\Pi(k)$（黏性项给出耗散），
两边移项即得。
\end{proof}

\textbf{实测}：$N=64,\nu=0.01$、$t=0.3$ 处，$k=1,\dots,32$ 上左边与右边的
最大偏差 $1.7\times10^{-6}$（时间方向用中心差分）。

\section{Burgers 的 $k^{-2}$ 段：是激波谱，不是惯性区}

\begin{theorem}[黏性截断波数]
激波厚度是 $\delta\sim\nu/u_0$，对应的耗散截断波数是 $k_d\sim1/\sqrt{\nu t}$；
能谱在 $k\ll k_d$ 上继承间断的 $k^{-2}$，在 $k\gtrsim k_d$ 上指数衰减。
\end{theorem}

\begin{proof}
Burgers 激波剖面 $u=-u_0\tanh(u_0x/(2\nu))$ 的 Fourier 变换是
$\mathrm{csch}(\pi\nu k/(2u_0))$：$k\ll k_d$ 时 $\propto k^{-1}$（幅度谱），
$k\gg k_d$ 时指数衰减。能量谱因此 $\propto k^{-2}$。
\end{proof}

\textbf{实测}（$N=128,\nu=0.02,t=1.5$）：窗口 $[3,8]$（上界 $\min(1.5k_d,N/3)$，
$k_d=5.77$）上的斜率 $-2.1994$（$R^2=0.999983$）。
但能流峰值与耗散率之比只有 $0.85$、能流末端 $-1.7\times10^{-2}$（本该回落到 $0$）：
一维没有 $k^{-5/3}$ 的惯性区，"$\Pi(k)\approx\varepsilon$"这条判据在这里给出的是
\textbf{假的不守恒}。

\textbf{实测（欠分辨）}：$\nu=0.002$（同一网格）时能流末端 $-4.280$、收支残差 $1.605$、
能量从 $0.500$ 涨到 $0.840$。收支残差随 $\nu$ 单调变差：
$1.35\times10^{-3},6.51\times10^{-3},3.45\times10^{-2},8.30\times10^{-2}$
（$\nu=0.03,0.02,0.01,0.005$，$N=64,t=1.2$）。

\begin{remark}[去混叠在积分量上几乎看不见]
$N=128,\nu=0.02,t=1.5$：去混叠与朴素乘积的总能量相对差 $2\times10^{-4}$，
收支残差 $5.9\times10^{-3}$ vs $2.8\times10^{-3}$（朴素那个还更小），
Enstrophy 差 $3.5\%$（$5.154$ vs $4.979$），逐点 $\max|\Delta u|/|u|=3.9\times10^{-2}$。
欠分辨时才在能量账上现形（$\nu=0.005$：Enstrophy 差 $6.5$ 倍、收支残差 $0.58$ vs $1.19$）。
\end{remark}

\section{四个可视化在看什么}

\begin{enumerate}
  \item \texttt{modal\_decay}（$4.6.1$ 模态系数衰减）：主图 $|a_k|$ 与理论包络，
        副图\textbf{局部斜率}，右图三类光滑性的误差阶梯。它可以切换两种基、四种剖面，
        参数 $N$ 与 $\rho$。要看的不是"曲线陡不陡"，而是局部斜率是否一直在变
        （默认帧 $-0.737\to-3.942$）以及跳变那条是否不下降（$1.956\to1.802$）。
  \item \texttt{dealias\_rule}（$4.6.2$ 乘积与 $3/2$ 规则）：主图是全带参照乘积谱与
        当前模式的乘积谱（竖线标出 $2/3$ 带边缘），副图是
        $|W_{\text{naive}}-W_{\text{ref,full}}|$ 的折叠谱——混叠\textbf{落在哪些模态}
        一眼可见；扫描图四条线分别是混叠、截断代价、只补零、以及 $2/3+3/2$。
  \item \texttt{convergence\_race}（$4.6.3$ 谱 vs 差分的误差—$N$）：三条误差曲线加
        两条每加倍比值曲线。招牌是比值：谱方法 $5.55\times10^{-2}\to2.90\times10^{-8}$、
        二阶差分恒在 $0.25$ 附近。时间轴上的点误差面板显示"误差长什么样"。
  \item \texttt{burgers\_cascade}（$4.6.4$ Burgers 级联）：能谱（叠加
        $k^{-2}$ 参考线与 $k_d$ 标记）、能流 $\Pi(k)$、解曲线、以及能量收支残差的时间曲线。
        参数是 $N$、$\nu$、去混叠开关、初值类型；$\nu$ 的临界线在 $0.015$。
\end{enumerate}

% 本文件由 tools/build_book.py 从 dist/reports/assertions.json 生成，请勿手改。
\begin{longtable}{@{}p{0.06\textwidth}p{0.32\textwidth}p{0.22\textwidth}p{0.16\textwidth}@{}}
\caption{章节 4.6 数值断言（生成时间 2026-10-18T18:30:41-0400）}\\
\toprule
类别 & 断言 & 理论值 & 实测值 \\
\midrule
\endfirsthead
\toprule
类别 & 断言 & 理论值 & 实测值 \\
\midrule
\endhead
\bottomrule
\endlastfoot

B & 角点帧总数 & 4 个 viz $\times$ 各自角点 $\times$ 时间端点 & 56 \checkmark \\
B & 检查过的掩码对 & 每个 *\_mask 都要有对应的底数组，且掩码内只剩地板值 & 14 \checkmark \\
B & viz1 的 fit\_\allowbreak{}p 与内核同源 & modal\_decay\_fit 给出的 -7.1290 & -7.12896 \checkmark \\
B & viz2 的四项账与内核同源 & alias /\allowbreak{} truncation /\allowbreak{} pad\_only /\allowbreak{} dealias 四个量各自独立 & 0.0119569 \checkmark \\
B & viz3 的差分阶与 shift 与内核同源 & -2 /\allowbreak{} -4，shift =\allowbreak{} c$\cdot$t =\allowbreak{} 0.510 & -1.97043 \checkmark \\
B & viz4 的能谱斜率与能量账与内核同源 & 窗口 [3,\allowbreak{}10] 上的幂律拟合 & -3.2071 \checkmark \\
B & convergence\_\allowbreak{}race 的时间端点确实不同 & t=\allowbreak{}0 与 t=\allowbreak{}tMax 的解不能相同 & 3.99855 \checkmark \\
B & burgers\_\allowbreak{}cascade 的时间端点确实不同 & t=\allowbreak{}0 与 t=\allowbreak{}tMax 的解不能相同 & 0.992988 \checkmark \\
B & modal\_\allowbreak{}decay 的帧可复现且参数生效 & 同参同 t 逐位相同；角点参数必须改变数组或 meta & 15 \checkmark \\
B & dealias\_\allowbreak{}rule 的帧可复现且参数生效 & 同参同 t 逐位相同；角点参数必须改变数组或 meta & 13 \checkmark \\
B & convergence\_\allowbreak{}race 的帧可复现且参数生效 & 同参同 t 逐位相同；角点参数必须改变数组或 meta & 15 \checkmark \\
B & burgers\_\allowbreak{}cascade 的帧可复现且参数生效 & 同参同 t 逐位相同；角点参数必须改变数组或 meta & 16 \checkmark \\
B & 裁剪掩码与 meta 计数一致 & n\_clipped =\allowbreak{} count(\allowbreak{}u\_clipped) & 8 \checkmark \\
A & n=16 三角多项式谱微分相对误差 & $\le$ 1e-13(\allowbreak{}舍入级) & $4.659\times 10^{-15}$ \checkmark \\
A & n=64 三角多项式谱微分相对误差 & $\le$ 1e-13(\allowbreak{}舍入级) & $1.009\times 10^{-14}$ \checkmark \\
A & 两种 N 上的谱微分误差都在舍入级 & $\le$ 1e-13(\allowbreak{}误差是 u 量级，不是「很小」) & $1.009\times 10^{-14}$ \checkmark \\
A & 跳变剖面谱微分的幅度随 N 增长 & $\approx$ (\allowbreak{}N\_hi/\allowbreak{}N\_lo)(\allowbreak{}Gibbs 层导数 \textasciitilde{} N) & 4.14636 \checkmark \\
A & rho=0.3 的系数偏离 rho\textasciicircum{}|k| & $\le$ 1e-14 & $5.551\times 10^{-17}$ \checkmark \\
A & rho=0.6 的系数偏离 rho\textasciicircum{}|k| & $\le$ 1e-06 & $2.245\times 10^{-11}$ \checkmark \\
A & rho=0.85 的混叠尾巴 & rho\textasciicircum{}(\allowbreak{}N/\allowbreak{}2) =\allowbreak{} 5.5e-3(\allowbreak{}采样网格的固有偏差) & 0.000411632 \checkmark \\
A & 朴素乘积在 k=-N/3 处的假模态 & 0.25(\allowbreak{}=\allowbreak{} cos\textasciicircum{}2 的 2K 分量折回；真值 0) & 0.25 \checkmark \\
A & n=16: D@x = 1 & $\le$ 1e-12 & $3.553\times 10^{-15}$ \checkmark \\
A & n=16: D@x\textasciicircum{}2 = 2x & $\le$ 1e-12 & $1.421\times 10^{-14}$ \checkmark \\
A & n=32: D@x = 1 & $\le$ 1e-12 & $5.684\times 10^{-14}$ \checkmark \\
A & n=32: D@x\textasciicircum{}2 = 2x & $\le$ 1e-12 & $7.105\times 10^{-14}$ \checkmark \\
A & |x| 的 Chebyshev 导数误差不随 n 下降 & O(\allowbreak{}1)：实测约 2.39(\allowbreak{}与 n 无关，跳变量为 2) & 2.38805 \checkmark \\
A & Gibbs 过冲随模态数的走向 & (\allowbreak{}2/\allowbreak{}pi)Si(\allowbreak{}pi)-1 =\allowbreak{} 0.1789797 & 0.000588555 \checkmark \\
A & dt=0.0025 与 DOP853 的相对偏差 & $\le$ 1e-10 & $1.426\times 10^{-12}$ \checkmark \\
A & 时间步从 0.04 缩到 0.0025 的偏差下降 & $\ge$ 1e4(\allowbreak{}四阶：16\textasciicircum{}4) & 65228.9 \checkmark \\
A & jump 的衰减指数 & -1 & -0.992147 \checkmark \\
A & corner 的衰减指数 & -2 & -1.98404 \checkmark \\
A & bump 的衰减指数 & -4 & -3.99992 \checkmark \\
A & 解析剖面的局部斜率一路变陡(不存在有限代数指数) & 每段至少再陡 2(\allowbreak{}指数衰减的可执行判据) & 3.58018 \checkmark \\
A & 尖角的逼近误差指数 & +1(\allowbreak{}a\_k\textasciitilde{}k\textasciicircum{}-2 的尾巴积分给出 N\textasciicircum{}-1) & 1.00376 \checkmark \\
A & 跳变的逼近误差不随 N 下降 & $\le$ 1.2(\allowbreak{}Gibbs 把误差钉住) & 1.05856 \checkmark \\
A & 解析剖面 N=256 的误差 & $\le$ 1e-13(\allowbreak{}地板) & $7.105\times 10^{-15}$ \checkmark \\
B & 合成-分解往返 & $\le$ 1e-13(\allowbreak{}舍入级) & $1.199\times 10^{-14}$ \checkmark \\
B & Parseval：sum\_\allowbreak{}k E\_\allowbreak{}k = mean(u\textasciicircum{}2) & $\le$ 1e-14 & $2.776\times 10^{-16}$ \checkmark \\
B & 带限输入的去混叠乘积 vs 高分辨率参照 & $\le$ 1e-13 & $1.035\times 10^{-16}$ \checkmark \\
B & 带限输入的朴素乘积偏差(保留带内) & $\le$ 1e-13(\allowbreak{}折返落在保留带外) & $1.438\times 10^{-16}$ \checkmark \\
B & 带限输入的朴素乘积偏差(保留带外) & 非零：折返就落在这里 & 0.00322086 \checkmark \\
B & 全带场的能量残差(quad 相位) & 朴素乘积应当非零，但量级随相位剧烈变化 & $7.654\times 10^{-18}$ \checkmark \\
B & 全带场的能量残差(lin 相位) & 朴素乘积应当非零，但量级随相位剧烈变化 & 0.00405371 \checkmark \\
B & 全带场的能量残差(rand 相位) & 朴素乘积应当非零，但量级随相位剧烈变化 & 0.00716665 \checkmark \\
B & 同谱不同相位的残差跨度 & $\ge$ 1e3(\allowbreak{}相位决定混叠是否打中能量账) & $9.364\times 10^{14}$ \checkmark \\
B & 带限场：朴素乘积的残差 & $\le$ 1e-15(\allowbreak{}折返在带外) & $3.424\times 10^{-18}$ \checkmark \\
B & 带限场：去混叠乘积的残差 & $\le$ 1e-15 & $8.446\times 10^{-17}$ \checkmark \\
B & 逐带预算的最大偏差(k=1..32) & $\le$ 1e-4(\allowbreak{}有限差分 + 账目) & $1.731\times 10^{-6}$ \checkmark \\
B & nu 从 0.03 降到 0.005 的收支残差 & 单调变差(\allowbreak{}欠分辨时 2/\allowbreak{}3 截断成为能量汇) & 0.0830384 \checkmark \\
B & c=1,t=1 的精确解 = 初值 & $\le$ 1e-14 & $4.885\times 10^{-15}$ \checkmark \\
B & c=1,t=1 时谱方法的误差 & $\le$ 1e-13(\allowbreak{}与 N 无关的退化) & $9.992\times 10^{-16}$ \checkmark \\
B & LGL 左端点 = +1 & $\le$ 1e-15 & 0 \checkmark \\
B & LGL 右端点 = -1 & $\le$ 1e-15 & 0 \checkmark \\
B & D 的行和 = 0(常数导数为 0) & $\le$ 1e-12 & $2.842\times 10^{-14}$ \checkmark \\
B & Runge(偶函数)的奇次系数 & $\le$ 1e-15 & $2.741\times 10^{-16}$ \checkmark \\
B & |x|(偶函数)的奇次系数 & $\le$ 1e-15 & $7.841\times 10^{-16}$ \checkmark \\
C & 二阶差分实测阶 & 2 & 1.97043 \checkmark \\
C & 四阶差分实测阶 & 4 & 3.96771 \checkmark \\
C & 二阶每加倍比值(最细两级) & 0.25(\allowbreak{}渐近值，粗网格端偏大) & 0.00476949 \checkmark \\
C & 四阶每加倍比值(最细两级) & 0.0625 & 0.00128826 \checkmark \\
C & 差分的每加倍比值从上方单调趋近理论值 & 单调下降且有下界 0.25 /\allowbreak{} 0.0625(\allowbreak{}粗网格端还没进渐近区) & 0.657009 \checkmark \\
C & 小 N 段四阶格式的假阶 & $\le$ 3.5(\allowbreak{}还没进渐近区) & 1.66885 \checkmark \\
C & 谱误差每加倍比值严格递减 & 严格递减(\allowbreak{}指数收敛) & $5.229\times 10^{-7}$ \checkmark \\
C & 谱方法每加倍比值 vs 差分每加倍比值 & $\le$ 0.01(\allowbreak{}差两个数量级以上) & $8.358\times 10^{-8}$ \checkmark \\
C & 谱误差的局部斜率逐段变陡 & 每段至少 +1.5(\allowbreak{}对数-对数图上向下弯) & 6.54561 \checkmark \\
C & 积分因子 RK4 实测时间阶 & 4 & 3.99315 \checkmark \\
C & 每减半步长的误差比值 & $\approx$ 1/\allowbreak{}16 =\allowbreak{} 0.0625 & 0.000740736 \checkmark \\
C & Runge 的误差每加倍比值递减 & 严格递减(\allowbreak{}指数) & $2.909\times 10^{-6}$ \checkmark \\
C & |x| 的 Chebyshev 误差指数 & +1(\allowbreak{}尖角 $\Rightarrow$ 误差 \textasciitilde{} n\textasciicircum{}-1) & 0.998951 \checkmark \\
C & sign(x) 的误差不随 n 下降 & $\le$ 1.2(\allowbreak{}跳变) & 1.02069 \checkmark \\
C & 能谱斜率(窗口 [3,8]) & -2(\allowbreak{}激波/\allowbreak{}间断的谱) & -2.19941 \checkmark \\
C & 欠分辨(nu=0.002)时能流末端 & 应当 $\approx$ 0(\allowbreak{}非线性项守恒总能量) & -4.28001 \checkmark \\
C & 截断代价 / 混叠误差 之比 & 2 \textasciitilde{} 7(\allowbreak{}截断这一笔更大) & 2.10352 \checkmark \\
C & 混叠误差随 decay 的幂律指数 & 负指数(\allowbreak{}谱越陡混叠越小) & -4.96326 \checkmark \\
\end{longtable}


\section{它会怎样骗你（本章版）}

\begin{enumerate}
  \item \textbf{把"谱方法准"当成方法的属性}：三角多项式的谱微分误差是
        $4.66\times10^{-15}$（$N=16$）与 $1.01\times10^{-14}$（$N=64$）——与 $N$ 无关；
        同一个算子在锯齿上让 $\max|u'|=6.73,13.80,27.92$（$N=32,64,128$），
        随 $N$ 线性增长。
  \item \textbf{拿双对数拟合的 $p$ 当收敛阶}：默认帧窗口 $k\in[4,32]$ 给
        $p=-7.129,R^2=0.938$，而局部斜率从 $-0.737$ 走到 $-3.942$。
        判据是每加倍比值（$0.143\to4.08\times10^{-4}$），不是某个 $p$。
  \item \textbf{对间断数据期待指数收敛}：跳变的 $L^\infty$ 误差
        $1.9078\to1.8022$（$N=16\to256$），Gibbs 过冲极限 $0.1789797$ 不为零，
        加密网格只把振荡压窄。
  \item \textbf{不筛零点模态就拟合}：采样方波的偶次谐波是常数 $1/N=0.03125$，
        给 $p=-0.848,R^2=0.346$；Chebyshev 的 Runge 奇次系数恰好为 $0$，
        不筛会得到 $+56$ 的正斜率。
  \item \textbf{把混叠与截断混成一笔账}：$N=128$、$\mathrm{decay}=1$ 时混叠
        $1.196\times10^{-2}$、截断代价 $3.740\times10^{-2}$、两者总差 $4.034\times10^{-2}$；
        截断的代价比它消掉的混叠大 $3$ 倍（比值随 $\mathrm{decay}$ 在 $2.10\sim6.74$ 之间）。
  \item \textbf{以为 $3/2$ 必须配 $2/3$}：只补零到精确 $3N/2$、完全不截断时偏差
        $7.76\times10^{-5}$，残差全部集中在 $|k|=N/2$ 一条边界模态上；
        $2/3$ 截断的价值是"不补零也能消混叠"。
  \item \textbf{把能量恒等式残差当去混叠的证据}：带限场的朴素与去混叠残差都是
        $10^{-17}$ 量级；全带场同一张谱换相位能从 $7.65\times10^{-18}$ 跳到
        $7.17\times10^{-3}$（跨度 $9.4\times10^{14}$）。
  \item \textbf{在整周期时刻量空间误差}：$c=1,t=1$ 时三种格式的误差都 $\le10^{-13}$，
        与 $N$ 无关——这是位移的整周期性，不是格式变准了。
        必须取非平凡位移（本帧 $c=0.85$，$shift=0.510$）。
  \item \textbf{把窗口藏起来报斜率}：四阶差分在 $N=16/32/64$ 上拟出 $1.669$，
        在 $N=128/192/256$ 上是 $3.9677$；谱曲线撞地板后，把地板点一起吃进最小二乘
        （	exttt{lo=64, hi=256}）能给 $p=-12.05,R^2=0.80$，而帧里的上半窗口
        （	exttt{lo=128}）在 $t=0.6/1.0$ 给 $-1.60/-1.33$（$R^2=0.09/0.06$）——
        两个窗口都不可信。斜率必须与窗口、$R^2$、用点数一起报。
  \item \textbf{以为换到 Chebyshev 就能救非光滑数据}：$|x|$ 的导数误差在
        $n=16..128$ 上恒为 $2.386$（$O(1)$），$\mathrm{sign}$ 的插值误差
        $1.9944\to1.9540$ 不动。
  \item \textbf{把一维 $k^{-2}$ 读成湍流惯性区}：默认帧斜率 $-2.1994$
        （$R^2=0.999983$）很漂亮，但 $\Pi_{\max}/\varepsilon=0.85$、能流末端
        $-1.7\times10^{-2}$；$\nu=0.002$ 时能量从 $0.500$ 涨到 $0.840$、
        收支残差 $1.605$。
  \item \textbf{以为去混叠总能改善能量守恒}：分辨良好时两种算法的总能量差
        $2\times10^{-4}$、收支残差 $5.9\times10^{-3}$ vs $2.8\times10^{-3}$
        （朴素更小），差别在 Enstrophy（$3.5\%$）与逐点解（$3.9\times10^{-2}$）上。
  \item \textbf{把显示用的裁剪当成物理}：解超过 $\pm1.6$ 的位置被截断显示，
        帧里必须同时给出 \texttt{u\_clipped} 掩码与 \texttt{n\_clipped} 计数
        （$\nu=0.002,t=1.5$ 时 $8$ 个点），否则读者会把"平台"当成激波结构。
\end{enumerate}
